\section{Síntesis y ampliación}

\subsection{Repaso de los puntos clave}

Esta clase abordó en profundidad la teoría y práctica de la guía durante el proceso de inferencia desde tres perspectivas:

\begin{enumerate}
    \item \textbf{Unificación desde la perspectiva SDE/ODE}: La guía mediante DPS no solo es aplicable a la formulación discreta de DDPM, sino que también se aplica consistentemente en las formulaciones de tiempo continuo de SDE y ODE.
    \item \textbf{Reinterpretación desde la perspectiva de RL}: Mediante un marco de maximización de recompensas sujeto a restricciones de KL, se derivó una fórmula de guía equivalente a DPS. Esta perspectiva de RL revela las profundas conexiones entre la guía durante la inferencia y el RLHF.
    \item \textbf{Amplia gama de aplicaciones}: Mediante técnicas de sincronización en el espacio canónico/espacio de instancia, los modelos preentrenados pueden generar contenido que va más allá del dominio de entrenamiento, como panoramas, texturas 3D e imágenes ambiguas.
\end{enumerate}

\subsection{Resonancia con el escalado en tiempo de prueba de LLM}

Esta clase situó la guía durante la inferencia de modelos difusos dentro de una tendencia más amplia de IA: el \textit{test-time scaling}. El campo de los LLM ha demostrado el enorme potencial de gastar más cómputo durante la inferencia a través de O1/O3/DeepSeek R1. La guía durante la inferencia de modelos difusos es esencialmente también una técnica de escalado en tiempo de prueba; no aumenta el costo del entrenamiento, sino que mejora la calidad de la salida mediante cálculos adicionales durante la inferencia.

\subsection{Tabla de resumen: métodos, valor y riesgos}

\begin{table}[H]
\centering
\begin{tabular}{p{0.22\textwidth} p{0.21\textwidth} p{0.23\textwidth} p{0.22\textwidth}}
\toprule
Tema & Valor central & Riesgo principal & Recomendación práctica \\
\midrule
DPS / Guía de problemas inversos & Manejo de restricciones de observación sin reentrenamiento & Grandes errores de aproximación en etapas de alto ruido & Utilizar un mayor número de pasos y realizar una programación de guía \\
Guía de recompensas desde la perspectiva de RL & Unificar el problema dentro del marco verifier + reward & Hacking de reward, aproximación de valor rugosa & Realizar primero la verificación con el verifier antes de abordar guías más fuertes \\
Generación sincronizada & Extrapolación de modelos preentrenados en 2D a panoramas / 3D / ilusiones & Conflictos de restricciones entre múltiples instancias, propenso a distorsiones & Registrar explícitamente proyecciones y regiones superpuestas entre el espacio canónico/espacio de instancia \\
Escalado en tiempo de prueba & Intercambiar mayor calidad de inferencia por menor costo de entrenamiento & Aumento del gasto computacional y latencia más alta & Controlar los costos mediante best-of-$N$ o asignación adaptativa de recursos \\
\bottomrule
\end{tabular}
\caption{Visión general de los métodos de la Clase 12}
\end{table}

\subsection{Tres puntos de acción para investigadores}

\begin{enumerate}
    \item Establecer primero una línea base estable en un reward simple, diferenciable y evaluable, antes de expandirse a configuraciones con múltiples restricciones.
    \item Crear un repositorio de casos de fallo unificado para cada método, registrando especialmente los patrones de degradación bajo pesos de guía altos y número reducido de pasos.
    \item Estudiar por comparación la guía durante la inferencia de modelos difusos frente a los métodos de verifier / reranking / cómputo adaptativo de LLM, en lugar de considerarlos por separado.
\end{enumerate}

\subsection{Preguntas adicionales para profundizar}

\begin{itemize}
    \item ¿Es posible entrenar un verifier unificado que sirva simultáneamente para problemas inversos, tareas de edición y preferencias estéticas?
    \item ¿Podría permitir que el modelo decida dinámicamente qué muestras merecen más pasos de inferencia, logrando así un verdadero escalado adaptativo en tiempo de prueba?
    \item Para modalidades más complejas como 3D, video y trayectorias robóticas, ¿sigue siendo estable la sincronización entre el espacio canónico/espacio de instancia?
\end{itemize}

\subsection{División de responsabilidades con los métodos de control durante el entrenamiento}

Al colocar el contenido de esta clase dentro de una perspectiva más completa del sistema generativo, se obtiene un juicio muy práctico: la guía durante la inferencia no busca reemplazar a ControlNet, LoRA o modelos difusos condicionales, sino colaborar con ellos mediante una división de tareas.

\begin{table}[H]
\centering
\begin{tabular}{p{0.22\textwidth} p{0.24\textwidth} p{0.22\textwidth} p{0.2\textwidth}}
\toprule
Enfoque & Problemas en los que es más hábil & Costo principal & Escenarios más adecuados \\
\midrule
Control durante el entrenamiento & Tipos de condiciones fijos y frecuentes & Necesidad de datos, entrenamiento y mantenimiento de múltiples modelos & Entornos de producción estables, servicios de baja latencia \\
Guía durante la inferencia & Restricciones temporales, variables y combinadas & Muestreo más lento, ajuste de hiperparámetros más sensible & Exploración de investigación, demandas de cola larga, ensamblaje de restricciones complejas \\
Enfoque híbrido & Estructura gruesa durante el entrenamiento y refinamiento fino durante la inferencia & Diseño del sistema más complejo & Tareas que requieren simultáneamente generación de alta calidad y control flexible \\
\bottomrule
\end{tabular}
\caption{Límites de responsabilidad entre el control durante el entrenamiento y la guía durante la inferencia}
\end{table}

La revelación central de esta tabla es: si una condición aparece casi siempre, lo mejor es absorberla en la etapa de entrenamiento; si las condiciones son temporales, transversales a tareas o difíciles de recopilar datos supervisados, es más adecuado mantenerlas para resolverlas durante la inferencia.

\subsection{Lecturas adicionales}

\begin{itemize}
    \item \textbf{Artículo original de DPS}: Chung et al., ``Diffusion Posterior Sampling for General Noisy Inverse Problems,'' ICLR 2023.
    \item \textbf{Reseña sobre escalado en tiempo de prueba}: El ponente recomendó un artículo de revisión sobre tecnologías de escalado en tiempo de prueba, que abarca diversas técnicas en los campos de modelos difusos y modelos de flujo.
    \item \textbf{SyncDiffusion}: Lee et al., ``SyncDiffusion: Coherent Montage via Synchronized Joint Diffusions,'' NeurIPS 2023——Generación sincronizada de imágenes panorámicas.
    \item \textbf{Visual Anagrams}: Geng et al., ``Visual Anagrams: Generating Multi-View Optical Illusions with Diffusion Models,'' CVPR 2024——Generación de ilusiones ópticas multivista.
    \item \textbf{TEXTure}: Richardson et al., ``TEXTure: Text-Guided Texturing of 3D Meshes,'' SIGGRAPH 2023——Generación de texturas 3D utilizando modelos difusos en 2D.
    \item \textbf{RLHF para Difusión}: Fan et al., ``DPOK: Reinforcement Learning for Fine-tuning Text-to-Image Diffusion Models''——Introducción de la idea de RLHF en el entrenamiento de modelos difusos texto-imagen.
\end{itemize}

%% --- Fin del contenido --- %%

\end{document}
